\chapter{门罗币地址 (Monero Addresses)}
\label{chapter:addresses}

存储在区块链中的数字货币的所有权由所谓的"地址"控制。地址是用来接收资金的，只有地址的持有者才能花费这些资金。\footnote{概率可忽略不计。}

更具体地说，一个地址拥有某些交易的"output"，这些output类似于票据，赋予地址持有者花费一定"金额"资金的权利。这样的一张票据可能写着"地址 C 现在拥有 5.3 XMR"。

要花费一个已拥有的 output，地址持有者将它作为新交易的 input 引用。这个新交易的 output 由其他地址拥有（或者如果发送者愿意，也可以由发送者自己的地址拥有）。交易的 input 总金额等于其 output 总金额，而且一旦被花费，已拥有的 output 就不能被再次花费。拥有地址 C 的 Carol 可以在新交易中引用那个旧的 output（例如"在这个交易中，我想花费那个旧的 output"），并添加一张票据说"地址 B 现在拥有 5.3 XMR"。

一个地址的余额是其未花费 output 中所含金额的总和。\footnote{被称为"钱包"的计算机应用程序用于查找和组织地址拥有的 output，保管其私钥以授权新交易，并将这些交易提交给网络进行验证和纳入区块链。}

我们将在第 \ref{chapter:pedersen-commitments} 章讨论如何对观察者隐藏金额，在第 \ref{chapter:transactions} 章讨论交易的结构（包括如何证明你正在花费一个已拥有的、之前未花费的 output，而不暴露具体花费的是哪个 output），在第 \ref{chapter:blockchain} 章讨论货币创造过程和观察者的角色。



\section{用户密钥 (User keys)}
\label{sec:user-keys}

与比特币不同，门罗币用户有两组私钥/公钥，\((k^v, K^v)\) 和 \( (k^s, K^s) \)，生成方式如第 \ref{ec:keys} 节所述。

用户的{地址}\marginnote{src/ common/ base58.cpp {\tt encode\_ addr()}}就是公钥对 \((K^v, K^s)\)。她的私钥就是对应的私钥对 \( (k^v, k^s) \)。\footnote{为了将地址传达给其他用户，最常见的方式是将其编码为 base-58，这是一种最初为比特币创建的二进制到文本编码方案 \cite{base-58-encoding}。更多详情请参见 \cite{luigi-address}。}\\

使用两组密钥可以实现功能隔离。其原理将在本章后面变得清晰，但现在让我们称私钥 $k^v$ 为{查看密钥 (view key)}，称 $k^s$ 为{花费密钥 (spend key)}。一个人可以使用他们的查看密钥来判断自己的地址是否拥有某个 output，而花费密钥则允许他们在交易中花费该 output（并追溯地确定它已被花费）。\footnote{它\marginnote{src/wallet/ api/wallet.cpp {\tt create()} wallet2.cpp {\tt generate()} {\tt get\_seed()}}目前最常见的做法是让查看密钥 $k^v$ 等于 $\mathcal{H}_n(k^s)$。这意味着一个人只需要保存他们的花费密钥 $k^s$，就可以访问（查看和花费）他们拥有的所有 output（已花费和未花费的）。花费密钥通常表示为 25 个单词的序列（其中第 25 个单词是校验和）。其他不太流行的方法包括：将 $k^v$ 和 $k^s$ 作为独立的随机数生成，或者生成一个随机的 12 助记词 $a$，其中 $k^s = {\tt sc\textunderscore reduce32}(\mathit{Keccak}(a))$，$k^v = {\tt sc\textunderscore reduce32}(\mathit{Keccak}(\mathit{Keccak}(a)))$。\cite{luigi-address}}



\section{一次性地址 (One-time addresses)}
\label{sec:one-time-addresses}

要接收资金，门罗币用户可以将自己的地址分发给其他用户，然后其他用户就可以通过交易的 output 向该地址发送资金。

地址永远不会被直接使用。\footnote{此处描述的方法不是由协议强制执行的，而是由钱包实现标准强制执行的。这意味着另一个钱包可以遵循比特币的风格，将接收者的地址直接包含在交易数据中。这样一个不合规的钱包将产生其他钱包无法使用的交易 output，而且由于 key image 的唯一性，每个类比特币风格的地址只能使用一次。}取而代之的是，对其应用类似于 Diffie-Hellman 的交换，为将要支付给用户的每个交易 output 创建一个唯一的{一次性地址 (one-time address)}。这样一来，即使知道所有用户地址的外部观察者也无法用它们来识别哪个用户拥有任何给定的交易 output。\footnote{概率可忽略不计。}

%接收者可以通过使用一次性地址签名消息来花费他们收到的output，从而证明他们知道私钥，因此拥有他们所花费的内容。我们将逐步充实这个概念。
%\\

让我们从一个非常简单的模拟交易开始，它只包含一个 output——一笔从 Alice 到 Bob 的"0"金额支付。

Bob 拥有私钥/公钥 $(k_B^v, k_B^s)$ 和 $(K_B^v, K_B^s)$，Alice 知道他的公钥（他的地址）。模拟交易可以如下进行 \cite{cryptoNoteWhitePaper}：

\begin{enumerate}
	\item Alice 生成一个随机数 $r \in_R \mathbb{Z}_l$，并计算一次性地址\footnote{在\marginnote{src/crypto/ crypto.cpp {\tt generate\_ key\_deri- vation()}}门罗币中，$r k^v G$ 的每个实例（包括当它用于交易的其他部分时）都乘以辅因子 8，所以在这种情况下 Alice 计算 $8*r K^v_B$，Bob 计算 $8*k^v_B r G$。据我们所知，这并没有什么实际用途（但它{确实}是用户必须遵守的规则）。乘以辅因子确保结果点在 $G$ 的子群中，但如果 $R$ 和 $K^v$ 本身就不共享子群，那么无论怎样，离散对数 $r$ 和 $k^v$ 都不能用来生成共享密钥。}\vspace{.175cm}
	\[K^o  = \mathcal{H}_n(r K_B^v)G + K_B^s\marginnote{src/crypto/ crypto.cpp {\tt derive\_pu- blic\_key()}}\]

	\item Alice 将 $K^o$ 设为收款地址，将 output 金额"0"和值 $r G$ 添加到交易数据中，并提交给网络。

	\item Bob\marginnote{src/crypto/ crypto.cpp {\tt derive\_ subaddress\_ public\_ key()}}接收到数据，看到值 $r G$ 和 $K^o$。他可以计算 $k_B^v r G = r K_B^v$。然后他可以计算 $K'^s_B = K^o - \mathcal{H}_n(r K_B^v)G$。当他看到 $K'^s_B = K_B^s$ 时，他就知道这个 output 是发给他的。\footnote{$K'^s_B $ 是通过 {\tt derive\_subaddress\_public\_key()} 计算的，因为普通地址的花费密钥存储在花费密钥查找表的索引 0 处，而子地址存储在索引 1、2……处。这很快就会有意义：见第 \ref{sec:subaddresses} 节。}

	私钥 $k_B^v$ 被称为"查看密钥"，因为任何拥有它（以及 Bob 的公钥花费密钥 $K_B^s$）的人都可以为区块链（交易记录）中的每个交易 output 计算 $K'^s_B$，并"查看"哪些属于 Bob。

	\item 该 output 的一次性密钥为：\vspace{.175cm}
	\begin{align*}
		K^o &= \mathcal{H}_n(r K_B^v)G + K_B^s = (\mathcal{H}_n(r K_B^v) + k_B^s)G  \\ 
		k^o &= \mathcal{H}_n(r K_B^v) + k_B^s
	\end{align*}
\end{enumerate}

要在新交易中花费他的"0"金额 output [原文如此]，Bob 所需要做的就是通过使用一次性密钥 $K^o$ 签名消息来证明所有权。私钥 $k_B^s$ 是"花费密钥"，因为证明 output 所有权需要它；而 $k_B^v$ 是"查看密钥"，因为它可以用来找到 Bob 可花费的 output。

正如将在第 \ref{chapter:transactions} 章中变得清楚的那样，没有 $k^o$ Alice 无法计算该 output 的 key image，因此她永远无法确定 Bob 是否花费了她发送给他的 output。\footnote{想象 Alice 产生了两个交易，每个都包含 Bob 可以花费的同一个一次性 output 地址 $K^o$。由于 $K^o$ 仅依赖于 $r$ 和 $K_B^v$，她没有理由不能这样做。Bob 只能花费其中一个 output，因为每个一次性地址只有一个 key image，所以如果不小心的话 Alice 可能会欺骗他。她可以创建一个给 Bob 大量资金的交易 1，然后创建一个给 Bob 少量资金的交易 2。如果他花费了交易 2 中的资金，他就永远无法花费交易 1 中的资金。事实上，没有人能花费交易 1 中的资金，实际上就"烧毁"了它。门罗币钱包已被设计为在此场景中忽略较小金额。}

Bob 可以将他的查看密钥交给第三方。这样的第三方可以是受信任的托管人、审计员、税务机关等。可以被允许读取用户交易历史而没有其他权限的人。这个第三方也将能够解密 output 的金额（将在第 \ref{sec:pedersen_monero} 节中解释）。有关 Bob 提供其交易历史信息的其他方式，请参见第 \ref{chapter:tx-knowledge-proofs} 章。


\subsection{多 output 交易 (Multi-output transactions)}
\label{sec:multi_out_transactions}

大多数交易将包含多个 output。至少，要将"找零"转回给发送者。\footnote{实际上，从协议 v12 起，每笔（非矿工）交易{必须}有两个 output，即使这意味着一个 output 的金额为 0（{\tt HF\_VERSION\_MIN\_2\_OUTPUTS}）。这通过将 1-output 情况与更常见的 2-output 交易混合来提高交易的不可区分性。在 v12 之前，核心钱包软件已经在创建 0 值的 output。核心实现将 0 金额的 output 发送到随机地址。\marginnote{src/wallet/ wallet2.cpp {\tt transfer\_ selected\_ rct()}}}\footnote{在 Bulletproofs 于协议 v8 中实现后，每笔交易被限制为不超过 16 个 output（{\tt BULLETPROOF\_MAX\_OUTPUTS}）。此前没有限制，只是对交易大小（字节数）有限制。}%BULLETPROOF_MAX_OUTPUTS; wallet2::transfer_selected_rct() for random address 0-amount

门罗币发送者通常只生成一个随机值 $r$。曲线点 $r G$ 通常被称为{交易公钥 (transaction public key)}，与其他交易数据一起发布在区块链中。\\

为了\marginnote{src/crypto- note\_core/ cryptonote\_ tx\_utils.cpp {\tt construct\_ tx\_with\_ tx\_key()}}[-.8cm]确保具有 $p$ 个 output 的交易中所有一次性地址都不同，即使同一收款地址被使用两次，门罗币使用了输出索引。交易的每个 output 都有一个索引 $t \in \{0, ..., p-1\}$。通过在 hash 之前将此值附加到共享密钥中，可以确保生成的一次性地址是唯一的：\marginnote{src/crypto/ crypto.cpp {\tt derive\_pu- blic\_key()}}[1.2cm]\vspace{.175cm}%construct_tx_with_tx_key() indexes 0 to p-1
\begin{align*}
  K_t^o &= \mathcal{H}_n(r K_t^v, t)G + K_t^s = (\mathcal{H}_n(r K_t^v, t) + k_t^s)G  \\ 
  k_t^o &= \mathcal{H}_n(r K_t^v, t) + k_t^s
\end{align*} 



\section{子地址 (Subaddresses)}
\label{sec:subaddresses}

门罗币用户可以从每个地址生成子地址 \cite{MRL-0006-subaddresses}。发送到子地址的资金可以使用其主地址的查看密钥和花费密钥来查看和花费。打个比方：一个在线银行账户可能有多个对应于信用卡和存款的余额，但它们都可以从同一个视角——账户持有人——访问和花费。\\

子地址对于在用户不想通过发布/使用同一地址来关联其活动时，将资金接收到同一地点非常方便。正如我们将看到的，大多数观察者将不得不解决离散对数问题才能确定一个给定的子地址是从任何特定地址派生的 \cite{MRL-0006-subaddresses}。\footnote{在子地址之前（在与协议 v7 对应的软件更新中添加 \cite{subaddress-pull-request}），门罗币用户可以简单地生成许多普通地址。要查看每个地址的余额，你需要对区块链记录进行单独扫描。这非常低效。有了子地址，用户维护一个（已 hash 的）花费密钥查找表，因此对 1 个子地址或 10,000 个子地址，扫描区块链所需的时间是相同的。}

它们对于区分接收到的 output 也很有用。例如，如果 Alice 想在周二从 Bob 那里买一个苹果，Bob 可以写一张描述购买的收据并为该收据创建一个子地址，然后让 Alice 在向他发送资金时使用该子地址。这样 Bob 就可以将他收到的钱与他卖出的苹果关联起来。我们将在下一节中探索另一种区分接收到的 output 的方法。%在子地址于2018年4月实现之前，门罗币开发者支持一种叫做支付ID的方法，它允许接收者区分拥有的output。支付ID是包含在交易'extra'、杂项数据中的文本字符串。接收者可以让发送者在写入交易时包含支付ID。支付ID也可以与所谓的集成地址一起编码。支付ID和集成地址不影响交易验证，所以任何人都可以实现它们，但由于它们目前在最流行的钱包中已被弃用，我们选择不在这里解释它们。}

Bob\marginnote{src/device/ device\_de- fault.cpp {\tt get\_sub- address\_ secret\_ key()}}从他的地址生成第 $i$ 个子地址 $(i = 1, 2, ...)$ 作为公钥对 $(K^{v,i}, K^{s,i})$：\footnote{实际上子地址的 hash 是域分离的，所以实际上是 $\mathcal{H}_n(T_{sub},k^v,i)$，其中 $T_{sub} = $"SubAddr"。为了简洁起见，我们在本文档中省略 $T_{sub}$。}\vspace{.175cm}
\begin{align*}
    K^{s,i} &= K^s + \mathcal{H}_n(k^v, i) G\\
    K^{v,i} &= k^v K^{s,i}
\end{align*}
\quad 所以，
\begin{alignat*}{2}
    K^{v,i} &= k^v&&(k^s + \mathcal{H}_n(k^v, i))G\\
    K^{s,i} &= &&(k^s + \mathcal{H}_n(k^v, i))G
\end{alignat*}
    

\subsection{向子地址发送 (Sending to a subaddress)}
    
假设 Alice 再次向 Bob 发送"0"金额，这次是通过他的子地址 $(K_B^{v,1}, K_B^{s,1})$。
\begin{enumerate}
	\item Alice 生成一个随机数 $r \in_R \mathbb{Z}_l$，并计算一次性地址\vspace{.175cm}
	\[ K^o  = \mathcal{H}_n(r K_B^{v,1},0)G + K_B^{s,1} \]

	\item Alice 将 $K^o$ 设为收款地址，将 output 金额"0"和值 $r K_B^{s,1}$ 添加到交易数据中，并提交给网络。
	
	\item Bob\marginnote{src/crypto/ crypto.cpp {\tt derive\_ subaddress\_ public\_ key()}}接收到数据，看到值 $r K_B^{s,1}$ 和 $K^o$。他可以计算 $k_B^v r K_B^{s,1} = r K_B^{v,1}$。然后他可以计算 $K'^{s}_B = K^o - \mathcal{H}_n(r K_B^{v,1},0)G$。当他看到 $K'^{s}_B = K^{s,1}_B$ 时，他就知道该交易是发给他的。\footnote{高级攻击者可能能够关联子地址 \cite{janus-attack}（即 Janus 攻击）。攻击者认为可能相关的子地址 $K_B^1$ $和$ $K_B^2$（其中一个可以是普通地址），他创建一个交易 output，其中 $K^o = \mathcal{H}_n(r K_B^{v,2},0)G + K_B^{s,1}$，并包含交易公钥 $r K_B^{s,2}$。Bob 计算 $r K_B^{v,2}$ 来找到 $K'^{s,1}_B$，但无法知道使用的是他的{另一个}（子）地址的密钥！如果他告诉攻击者他收到了发往 $K_B^1$ 的资金，攻击者就会知道 $K_B^2$ 是一个相关的子地址（或普通地址）。由于子地址不在协议范围内，缓解措施取决于钱包实现者。目前没有已知的钱包这样做过，而且任何缓解措施只对合规钱包有效。潜在的缓解措施包括：不向攻击者通报收到的资金、在交易数据中包含加密的交易私钥 $r$、使用 $K^{s,1}$ 而不是 $G$ 作为基点对共享密钥进行 Schnorr 签名、在交易数据中包含 $rG$ 并使用 $rK^{s,1} \stackrel{?}{=} (k^s + \mathcal{H}_n(k^v, 1))*rG$ 验证共享密钥（需要私有花费密钥）。普通地址收到的 output 也应该被验证。有关最后一种缓解措施的讨论，请参见 \cite{janus-mitigation-issue-62}。}
	
	Bob 只需要他的私有查看密钥 $k_B^v$ 和子地址的公钥花费密钥 $K^{s,1}_B$ 就可以找到发送到他子地址的交易 output。
	
	\item 该 output 的一次性密钥为：\vspace{.175cm}
	\begin{align*}
		K^o &= \mathcal{H}_n(r K_B^{v,1},0)G + K_B^{s,1} = (\mathcal{H}_n(r K_B^{v,1},0) + k_B^{s,1})G  \\ 
		k^o &= \mathcal{H}_n(r K_B^{v,1},0) + k_B^{s,1}
	\end{align*}
\end{enumerate}

现在，\marginnote{src/crypto- note\_core/ cryptonote\_ tx\_utils.cpp {\tt construct\_ tx\_with\_ tx\_key()}}Alice 的交易公钥是 Bob 特有的（$r K_B^{s,1}$ 而不是 $r G$）。如果她创建一个包含 $p$ 个 output 的交易，且至少有一个 output 是发往子地址的，Alice 需要为每个 output $t \in \{0,...,p-1\}$ 创建一个唯一的交易公钥。换句话说，如果 Alice 向 Bob 的子地址 $(K_B^{v,1}, K_B^{s,1})$ 和 Carol 的地址 $(K_C^v, K_C^s)$ 发送，她将在交易数据中放置两个交易公钥 \{$r_1 K_B^{s,1},r_2 G$\}。\footnote{在\marginnote{src/crypto- note\_basic/ cryptonote\_ basic\_impl.cpp {\tt get\_account\_ address\_as\_ str()}}门罗币中，子地址以"8"为前缀，与以"4"为前缀的地址区分开来。这有助于发送者在构建交易时选择正确的过程。}\footnote{关于何时使用额外交易公钥有一些复杂性（参见代码路径 {\tt transfer\_selected\_rct()} $\rightarrow$ {\tt construct\_tx\_and\_get\_tx\_key()} $\rightarrow$ {\tt construct\_tx\_with\_tx\_key()} $\rightarrow$ {\tt generate\_output\_ephemeral\_keys()} 和 {\tt classify\_addresses()}）围绕找零 output，交易作者知道接收者的查看密钥（因为是他自己；当需要 0 金额 output 时创建的虚拟找零 output 也是如此，因为作者生成这些地址）。只要至少有两个非找零 output，并且其中至少一个的接收者是子地址，就按照上面解释的正常方式进行（核心实现中的一个当前 bug 会在额外密钥之外再向交易数据添加一个额外的交易公钥，但该公钥不用于任何目的）。如果只有找零 output 是发往子地址的，或者只有一个非找零 output 且它是发往子地址的，那么只创建一个交易公钥。在前一种情况下，交易公钥是 $...

\section{集成地址 (Integrated addresses)}
\label{sec:integrated-addresses}

为了区分接收到的 output，接收者可以请求发送者在交易数据中包含一个{支付 ID (payment ID)}。\footnote{目前，门罗币实现每笔交易只支持一个支付 ID，无论它有多少个 output。}例如，如果 Alice 想在周二从 Bob 那里买一个苹果，Bob 可以写一张描述购买的收据，并让 Alice 在向他发送资金时包含收据的 ID 号。这样 Bob 就可以将他收到的钱与他卖出的苹果关联起来。

曾经发送者可以明文传送支付 ID，但在交易中手动包含 ID 不太方便，而且明文对接收者来说是隐私隐患，可能会无意中暴露他们的活动。\footnote{这些长格式（256 位）明文支付 ID 在核心软件实现的 v0.15 中已被弃用（与 2019 年 11 月的协议 v12 同时）。虽然其他钱包可能仍然支持它们并允许它们包含在交易数据中，但核心钱包将忽略它们。}在\marginnote{src/crypto- note\_basic/ cryptonote\_ basic\_ impl.cpp {\tt get\_ account\_ integrated\_ address\_as\_ str()}}门罗币中，接收者可以将支付 ID 集成到他们的地址中，并将这些{集成地址}（包含 $K^v$、$K^s$、支付 ID）提供给发送者。支付 ID 在技术上可以集成到任何类型的地址中，包括普通地址、子地址和多签地址。\footnote{据作者所知，集成地址只对普通地址实现过。}

发送者向集成地址发送 output 时，可以使用共享密钥 $r K_t^v$ 和 XOR 操作（回顾第 \ref{sec:XOR_section} 节）对支付 ID 进行编码，然后接收者可以使用相应的交易公钥和另一次 XOR 操作进行解码 \cite{integrated-addresses}。以这种方式编码支付 ID 还允许发送者证明他们制作了特定的交易 output（例如用于审计、退款等）。\footnote{由于观察者可以识别有支付 ID 和没有支付 ID 的交易之间的差异，使用它们被认为会使门罗币交易历史不那么统一。因此，从协议 v10 起，核心实现向所有 2-output 交易添加一个虚拟的\marginnote{src/crypto- note\_core/ cryptonote\_ tx\_utils.cpp {\tt construct\_ tx\_with\_ tx\_key()}}加密支付 ID。解码一个会显示全是 0（这不是协议要求的，只是最佳实践）。}


\subsubsection*{编码 (Encoding)}

发送者\marginnote{src/device/ device\_de- fault.cpp {\tt encrypt\_ payment\_ id()}}编码支付 ID 以包含在交易数据中\footnote{在门罗币中，集成地址的支付 ID 通常为 64 位长。}\vspace{.175cm}
\begin{align*}
         k_{\textrm{mask}} &= \mathcal{H}_n(r K_t^v,\textrm{pid\_tag}) \\
      k_{\textrm{payment ID}} &= k_{\textrm{mask}} \rightarrow \textrm{reduced to bit length of payment ID}\\
  \textrm{encoded payment ID} &= \textrm{XOR}(k_{\textrm{payment ID}}, \textrm{payment ID})
\end{align*}

我们包含 pid\_tag 以确保 $k_{\textrm{mask}}$ 与一次性 output 地址中的分量 $\mathcal{H}_n(r K_t^v, t)$ 不同。\footnote{在门罗币中，pid\_tag = {\tt ENCRYPTED\_PAYMENT\_ID\_TAIL} = 141。在例如多 input 交易中，我们计算 $\mathcal{H}_n(r K_t^v, t) \pmod l$ 以确保我们使用的是小于 EC 子群阶的标量，但由于 $l$ 是 253 位而支付 ID 只有 64 位，对编码支付 ID 取模将没有意义，所以我们不这样做。}


\subsubsection*{解码 (Decoding)}

支付 ID 是为哪个接收者创建的，该接收者\marginnote{src/device/ device.hpp {\tt decrypt\_ payment\_ id()}}可以使用他的查看密钥和交易公钥 $r G$ 来找到它：\footnote{交易数据不指示支付 ID"属于"哪个 output。接收者必须识别自己的支付 ID。}\vspace{.175cm}
\begin{align*}
         k_{\textrm{mask}} &= \mathcal{H}_n(k_t^v r G,\textrm{pid\_tag}) \\
      k_{\textrm{payment ID}} &= k_{\textrm{mask}} \rightarrow \textrm{reduced to bit length of payment ID}\\
          \textrm{payment ID} &= \textrm{XOR}(k_{\textrm{payment ID}}, \textrm{encoded payment ID})
\end{align*}

同样，发送者可以通过重新计算共享密钥来解码他们之前编码的支付 ID。%删掉这行？


\section{多签地址 (Multisignature addresses)}
\label{sec:multisignature-addresses}

有时在不同的人/地址之间共享资金的所有权是有用的。我们将在第 \ref{chapter:multisignatures} 章详细讨论这个话题。